\documentclass[11pt,a4paper]{article}

\usepackage[margin=2.4cm]{geometry}
\usepackage{amsmath,amssymb,amsfonts}
\usepackage{booktabs}
\usepackage{longtable}
\usepackage{array}
\usepackage{siunitx}
\usepackage[version=4]{mhchem}
\usepackage{xcolor}
\usepackage[colorlinks=true,linkcolor=blue!60!black,citecolor=teal,urlcolor=blue!60!black]{hyperref}
\usepackage{enumitem}
\usepackage{mathtools}
\usepackage{fancyhdr}

\newcommand{\dd}[2]{\frac{\mathrm{d}#1}{\mathrm{d}#2}}
\newcommand{\pd}[2]{\frac{\partial #1}{\partial #2}}
\newcommand{\Tcan}{T_{1}}
\newcommand{\Tair}{T_{2}}
\newcommand{\Cair}{C_{1}}
\newcommand{\Vair}{V_{1}}
\newcommand{\Tflo}{T_{5}}

\pagestyle{fancy}
\fancyhf{}
\rhead{Greenhouse--Cucumber Digital Twin}
\lhead{Model description}
\cfoot{\thepage}

\title{\textbf{A Mechanistic Greenhouse--Cucumber Digital Twin}\\[4pt]
\large Climate, Photosynthesis and Crop--Growth Subsystems:\\ Governing Equations, Variables and Control}
\author{Greenhouses Project --- Model Reference}
\date{\today}

\begin{document}
\maketitle

\begin{abstract}
\noindent This document is a complete, self-contained specification of the mechanistic
greenhouse--cucumber model that constitutes the \emph{digital twin}. The model couples three
subsystems on two time scales: a fast \textbf{climate} subsystem (a system of ordinary
differential equations for the greenhouse micro-climate, integrated within each day), a
\textbf{photosynthesis} subsystem (a Farquhar--von~Caemmerer--Berry biochemical leaf model scaled
to the canopy), and a daily \textbf{crop-growth} subsystem (a source--sink carbon-allocation map
with organ pools, thermal-time development, carbohydrate-regulated flowering and fruit cohorts).
Every differential equation is given with its complete right-hand side; every state variable,
input and control is defined where it first appears; and the physical, biochemical and
physiological principles behind each block are explained. Numerical values of all constants are
collected in the appendices. Parameter-inference and reinforcement-learning layers built on top
of this model are intentionally \emph{out of scope} here.
\end{abstract}

\tableofcontents
\bigskip

%======================================================================
\section{Overview and modelling philosophy}\label{sec:overview}
%======================================================================

\paragraph{What this section contains.} A bird's-eye view of the model: the three subsystems,
the variables they exchange, the two-time-scale coupling, and the notation used throughout.

\subsection{The three subsystems and their coupling}
The digital twin predicts how a greenhouse cucumber crop develops under a given weather series and
a given control (setpoint) programme. It is organised as three coupled subsystems:

\begin{enumerate}[leftmargin=2.2em]
  \item \textbf{Climate} (Section~\ref{sec:climate}). A lumped-parameter energy- and mass-balance
  model of the greenhouse compartment. Its state is the vector
  \begin{equation}
    \mathbf{x}_{\text{clim}} = \big(\,\Tcan,\ \Tair,\ \Cair,\ \Vair,\ \Tflo\,\big),
  \end{equation}
  the canopy temperature $\Tcan$, the greenhouse-air temperature $\Tair$, the air
  \ce{CO2} concentration $\Cair$, the air vapour pressure $\Vair$, and a floor/soil thermal-mass
  temperature $\Tflo$. These obey the ODEs of Section~\ref{sec:climate}, driven by outside weather
  and by the actuators.
  \item \textbf{Photosynthesis} (Section~\ref{sec:photo}). A biochemical model that converts the
  instantaneous canopy light, temperature and \ce{CO2} into a gross assimilation rate $A$. It is
  evaluated \emph{inside} the climate ODE (it both consumes \ce{CO2} in the air balance and,
  integrated over the day, feeds the crop).
  \item \textbf{Crop growth} (Section~\ref{sec:growth}). A daily map that spends the day's
  assimilate on maintenance respiration and organ growth, advances development in thermal time,
  sets and grows fruit, harvests mature fruit, and updates the leaf area index (LAI).
\end{enumerate}

\subsection{Two time scales}
The climate and photosynthesis subsystems live on a \emph{fast} (sub-daily) scale: within one day
the ODE of Section~\ref{sec:climate} is integrated over $[0,\SI{86400}{s}]$ at a \emph{fixed} LAI,
and the assimilation rate $A$ is accumulated,
\begin{equation}
  P_g \;=\; \Big(\!\int_{0}^{\SI{86400}{s}}\! A(t)\,\mathrm{d}t\Big)\cdot 10^{-3}\cdot\frac{30}{44}
  \qquad [\si{g\,CH_2O.m^{-2}.d^{-1}}],
  \label{eq:Pg}
\end{equation}
where the factor $10^{-3}$ converts \si{mg\to g} and $30/44$ converts assimilated \ce{CO2}
(molar mass \SI{44}{g/mol}) to carbohydrate \ce{CH2O} (molar mass \SI{30}{g/mol}). The crop-growth
subsystem lives on the \emph{slow} (daily) scale: once per day it receives $P_g$ and the daily-mean
air temperature $\overline{\Tair}$, updates the crop state, and returns the new LAI, which becomes
the fixed LAI for the next day's climate integration. This one-way-per-day hand-off (LAI $\to$
climate/photosynthesis $\to$ $P_g$ $\to$ crop $\to$ LAI) is the backbone of the twin.

\subsection{Notation}
Throughout, $\Tcan,\Tair,\Tflo$ are absolute temperatures in kelvin; the crop model works with the
daily-mean air temperature in \si{\degreeCelsius}. Symbols beginning with $I$ (with numeric
subscripts $I_1,\dots,I_{11}$) are \emph{inputs} to the climate right-hand side --- either weather
drivers or quantities supplied by the controller (Table~\ref{tab:inputs}); symbols
$U_1,\dots,U_{12}$ are the \emph{control actuators} (Table~\ref{tab:controls}). Greek-letter
constants ($\alpha_i,\beta_i,\gamma_i,\dots$) are fixed physical parameters listed in
Appendix~\ref{app:climate}. Intermediate flux quantities are named by a letter and index
($h_i$ sensible-heat fluxes, $l_i$ latent-heat fluxes, $r_i$ radiative fluxes, $f_i$ ventilation
exchange rates, $o_i$ \ce{CO2} fluxes, $p_i$ vapour fluxes, $q_i$ transpiration auxiliaries); they
are all defined in Section~\ref{sec:climate}.

%======================================================================
\section{Climate subsystem}\label{sec:climate}
%======================================================================

\paragraph{What this section contains.} The complete micro-climate model: the five state
equations (canopy temperature, air temperature, \ce{CO2}, vapour pressure and floor thermal mass),
followed by the submodels that define their right-hand-side terms --- shortwave (PAR/NIR) and
longwave (FIR) radiation, natural ventilation, canopy transpiration and its stomatal control,
condensation, and psychrometrics. The underlying principle throughout is the \textbf{lumped
energy and mass balance}: each compartment (canopy, air, floor) is treated as spatially uniform,
and its temperature/concentration changes at a rate equal to the net flux into it divided by its
heat capacity or the compartment volume. The formulation follows the Vanthoor greenhouse model.

\subsection{Governing state equations}

\paragraph{Canopy temperature $\Tcan$.} The canopy leaves store heat with capacity
$\alpha_1 I_1$ per unit ground area, where $I_1=\mathrm{LAI}$. Their temperature responds to
absorbed short-wave radiation ($r_1$, PAR, and $r_5$, NIR + lamp), long-wave exchange with the
heating pipe ($r_6$) and with sky/floor ($r_7,r_{14}$), sensible-heat convection to the air
($h_1$) and latent-heat loss by transpiration ($l_1$):
\begin{equation}
  \dd{\Tcan}{t} \;=\; \frac{1}{\alpha_1 I_1}\Big(\,r_1 + r_5 + r_6 - h_1 - l_1 - r_7 - r_{14}\,\Big).
  \label{eq:dT1}
\end{equation}

\paragraph{Greenhouse-air temperature $\Tair$.} The air (mean height $\phi_2$, density $\rho_3$,
specific heat $\alpha_5$) is heated by convection from the canopy ($h_1$), by the fan--pad/forced
inlet ($h_2$), by the heating pipe ($h_4$) and the lamps' direct/reflector heat ($h_{12}$), and by
absorbed radiation on building elements ($r_8$); it loses heat by ventilation to outside ($h_7$),
by conduction through the soil path ($h_{11}$), by the net long-wave balance ($r_{10}$) and by fog
evaporation ($l_2$). A separate floor-exchange term $\dot{T}_{2,\text{floor}}$
(Section~\ref{sec:floor}) couples it to the thermal mass:
\begin{equation}
  \dd{\Tair}{t} \;=\; \frac{1}{\phi_2\,\rho_3\,\alpha_5}
     \Big(\,h_1 + h_2 + h_4 - h_7 - h_{11} + h_{12} + r_8 - r_{10} - l_2\,\Big)
     \;+\; \dot{T}_{2,\text{floor}}.
  \label{eq:dT2}
\end{equation}

\paragraph{Air \ce{CO2} concentration $\Cair$.} Mass balance of \ce{CO2} (in \si{mg.m^{-3}}) over
the air column of height $\phi_2$: sources are the direct air-heater combustion ($o_1$) and the
external \ce{CO2} dosing valve ($o_2$); exchange with outside occurs through the fan--pad ($o_3$)
and through ventilation + leakage ($o_5$); the canopy removes \ce{CO2} by photosynthesis at the
assimilation rate $o_4=A$:
\begin{equation}
  \dd{\Cair}{t} \;=\; \frac{1}{\phi_2}\Big(\,o_1 + o_2 + o_3 - o_4 - o_5\,\Big).
  \label{eq:dC1}
\end{equation}

\paragraph{Air vapour pressure $\Vair$.} Water-vapour balance (in \si{Pa}): the canopy adds vapour
by transpiration ($p_1$), the fan--pad ($p_2$), the fog system ($p_3$) and the air heater ($p_4$);
vapour is removed by ventilation ($p_5$), the fan--pad outflow ($p_6$) and condensation on the
mechanical cooler ($p_7$), the thermal screen ($p_8$) and the cover ($p_9$). The capacity factor
$\kappa_3=\psi_1\phi_2/(\omega_2 \Tair)$ converts the vapour mass balance to a pressure rate through
the ideal-gas law:
\begin{equation}
  \dd{\Vair}{t} \;=\; \frac{1}{\kappa_3}
     \Big(\,p_1 + p_2 + p_3 + p_4 - p_5 - p_6 - p_7 - p_8 - p_9\,\Big).
  \label{eq:dV1}
\end{equation}

\paragraph{Floor / soil thermal mass $\Tflo$.} See Section~\ref{sec:floor}.

\medskip
The auxiliary flux terms in Eqs.~\eqref{eq:dT1}--\eqref{eq:dV1} are defined in
Sections~\ref{sec:rad}--\ref{sec:psychro}. We first list the transpiration auxiliaries $q_i$,
because several balances use $q_1$ (the vapour-exchange coefficient) and $q_2=\mathrm{Pws}(\Tcan)$
(saturated vapour pressure at leaf temperature).

\subsection{Transpiration coefficient and stomatal control}\label{sec:transp}
Canopy transpiration is a diffusion process: water vapour flows from the saturated interior of the
leaf (vapour pressure $q_2=\mathrm{Pws}(\Tcan)$) to the drier greenhouse air (vapour pressure
$\Vair$), across the sum of a stomatal resistance $q_3$ and a fixed boundary-layer resistance. The
coupling coefficient is
\begin{equation}
  q_1 \;=\; \frac{2\,\rho_3\,\alpha_5\,I_1}{\gamma\,\gamma_2\,(\gamma_3 + q_3)},
\end{equation}
so that the transpiration vapour flux is $q_1(q_2-\Vair)$ and the associated latent-heat loss is
$l_1=\gamma_2 q_1 (q_2-\Vair)$, with $\gamma_2$ the latent heat of vaporisation and $\gamma$ the
psychrometric constant. The \textbf{stomatal resistance} $q_3$ increases when the plant closes its
stomata --- in the dark, at high \ce{CO2}, and under large vapour-pressure deficit --- following the
Vanthoor switch functions:
\begin{align}
  q_7 &= \big(1 + e^{(I_9-\delta_1)/\gamma_5}\big)^{-1}, &
  q_8 &= \delta_4 + (\delta_5-\delta_4)\,q_7, \\
  q_9 &= \delta_6 + (\delta_7-\delta_6)\,q_7, &
  q_{10} &= \frac{I_9+\delta_2}{I_9+\delta_3}, \\
  q_4 &= 1 + q_8\,(\Cair-200)^2, &
  q_5 &= 1 + q_9\,(q_2-\Vair)^2, \\
  q_3 &= \gamma_4\, q_{10}\, q_4\, q_5. &&
\end{align}
Here $I_9$ is the PAR above the canopy (a radiation input, Section~\ref{sec:rad}); $q_7$ is a smooth
day/night switch centred at the light level $\delta_1$; $q_{10}$ is the radiation dependence of the
minimum resistance $\gamma_4$; $q_4$ and $q_5$ inflate the resistance far from
$\ce{CO2}=\SI{200}{ppm}$ and at large leaf-to-air vapour deficit, respectively. The quantity
$q_6=\mathrm{Pws}(I_6)$ (saturation pressure at the mechanical-cooler temperature $I_6$) is used by
the condensation terms.

\subsection{Radiation submodel}\label{sec:rad}
\paragraph{Principle.} Radiative transfer is split by wavelength band: \emph{PAR}
(photosynthetically active, $\SI{400}{}$--$\SI{700}{nm}$), \emph{NIR} (near-infrared, solar heat),
and \emph{FIR} (far-infrared, thermal emission). Short-wave (PAR+NIR) penetration into the canopy
follows the \textbf{Beer--Lambert law}: a beam is attenuated as $e^{-\beta\,\mathrm{LAI}}$ with an
extinction coefficient $\beta$ specific to the band. Long-wave exchange between surfaces follows the
\textbf{Stefan--Boltzmann law}, net flux $\propto \sigma(T_a^4-T_b^4)$, weighted by emissivities and
view factors.

\paragraph{Canopy fractions.} With $I_1=\mathrm{LAI}$,
\begin{align}
  g_1 &= 0.49\,\big(1-e^{-\beta_3 I_1}\big), &
  g_2 &= 1 - U_1(1-\tau_3), &
  g_3 &= \tau_1\,g_2\,\big(1-0.49\pi\gamma_1\phi_1\big)e^{-\beta_3 I_1},
\end{align}
where $g_1$ is the NIR-absorbing canopy fraction, $g_2$ the FIR transmission through the (possibly
deployed, $U_1$) thermal screen of transmissivity $\tau_3$, and $g_3$ a combined cover--screen--canopy
FIR factor.

\paragraph{PAR absorbed by the canopy.} The transmitted outside PAR reaching the top of the canopy is
$r_4=(1-\eta_1)\tau_1\eta_2 I_2$ ($I_2$ = outside global radiation, $\eta_2$ its PAR fraction,
$\tau_1$ cover PAR transmissivity, $\eta_1$ the building-absorbed fraction). Direct and
floor-reflected PAR absorbed by the canopy are
\begin{align}
  r_2 &= r_4\,(1-\rho_1)\big(1-e^{-\beta_1 I_1}\big), \\
  r_3 &= r_4\,e^{-\beta_1 I_1}\,\rho_2\,(1-\rho_1)\big(1-e^{-\beta_2 I_1}\big), \\
  r_1 &= r_2 + r_3,
\end{align}
with $\rho_1$ the canopy PAR reflectance, $\rho_2$ the floor PAR reflectance, and
$\beta_1,\beta_2$ the extinction coefficients for downward and floor-reflected PAR.

\paragraph{NIR and lamp short-wave to canopy and building.}
\begin{align}
  r_5 &= (1-\eta_1)\,\alpha_2\,\big(\eta_3 I_2 + \eta_{14}\,\alpha_{12}\,U_{12}\big), \\
  r_8 &= I_2\Big(\eta_1\big(\tau_1\eta_2 + (\alpha_2+\alpha_7)\eta_3\big) + \big(\alpha_8\eta_2 + \alpha_9\eta_3\big)\Big),
\end{align}
where $\alpha_2$ is the canopy NIR absorptance, $\eta_3$ the NIR fraction of global radiation,
$\alpha_{12}U_{12}$ the lamp radiation ($U_{12}$ the lamp on/off actuator) and $\eta_{14}$ its NIR
share; $r_8$ collects short-wave absorbed by building/floor/cover.

\paragraph{Long-wave (FIR) exchange.} With the sky-directed Stefan--Boltzmann kernel
$B_{\text{sky}}=\epsilon_3\sigma(\Tair^4-I_4^4)$ ($I_4$ = sky temperature, Section~\ref{sec:sky}),
\begin{align}
  r_{11} &= \epsilon_4\,g_3\,\epsilon_3\,\sigma\,(I_7^4 - I_4^4), &
  r_{12} &= \epsilon_{scr}\,(U_1\tau_2)\,B_{\text{sky}}, \\
  r_{13} &= \epsilon_6\,B_{\text{sky}}, &
  r_{10} &= r_{11}+r_{12}+r_{13},\\
  r_{14} &= \epsilon_4\,\epsilon_2\,\big(1-0.49\pi\gamma_1\phi_1\big)\sigma\,(\Tair^4 - I_7^4), &
  r_{15} &= \epsilon_2\,\epsilon_5\,\big(1-e^{-\beta_2 I_1}\big)g_2\,\sigma\,(\Tcan^4 - I_4^4).
\end{align}
$r_{10}$ is the net FIR loss of the air/cover system to the sky ($r_{13}$ through the cover of
emissivity $\epsilon_6$, $r_{12}$ through the deployed screen of \emph{energy-screen} emissivity
$\epsilon_{scr}$, $r_{11}$ the floor--sky path via $I_7$). $r_{14}$ is air$\leftrightarrow$floor FIR;
$r_{15}$ the canopy--sky FIR. The heating-pipe and canopy--sky canopy terms are
\begin{align}
  r_6 &= \alpha_3\,\epsilon_1\,\epsilon_2\,g_1\,\sigma\,(I_3^4 - \Tcan^4), &
  r_7 &= \frac{g_1}{0.49}\,\epsilon_2\,\epsilon_3\,g_2\,\sigma\,(\Tcan^4 - I_4^4),
\end{align}
with $I_3$ the heating-pipe temperature and $\alpha_3,\epsilon_1$ the pipe area and emissivity.

\subsection{Ventilation submodel}\label{sec:vent}
\paragraph{Principle.} Natural ventilation exchanges air with the outside through roof and side
vents. Two mechanisms drive it: \emph{wind} (dynamic pressure $\propto I_8^2$, $I_8$ = wind speed)
and \emph{buoyancy} (the stack effect, $\propto g\,\Delta T/\overline{T}$). A ventilation rate is a
non-negative air-exchange rate; the \emph{direction} of any heat/\ce{CO2}/vapour transport is carried
by the corresponding concentration or temperature difference in the balance, not by the rate.

\paragraph{Aperture-dependent rates.} With $U_6$ the side-vent and $U_8$ the roof-vent apertures,
\begin{align}
  n_1 &= \nu_1(1-\eta_{10}U_5), & n_2 &= \nu_3 U_6, & n_3 &= \nu_2(1-\eta_{11}U_5),\\
  f_5 &= \frac{n_1\,n_2\,I_8\,\sqrt{n_3}}{2\alpha_6}, &
  f_{7t} &= \max\!\Big(\omega_1\nu_6\frac{\Tair-I_5}{\Tair+I_5} + n_3 I_8^2,\ 0\Big), &
  f_7 &= \frac{U_8\,\nu_5\,n_1}{2\alpha_6}\sqrt{f_{7t}},
\end{align}
where $f_5$ is the wind-driven side-vent exchange, $f_7$ the roof-vent exchange combining buoyancy
(the $\omega_1\nu_6(\Tair-I_5)/(\Tair+I_5)$ stack term, $\omega_1=g$, $I_5$ = outside temperature)
and wind ($n_3 I_8^2$), and $\alpha_6$ the floor area. The infiltration/leakage rate has a wind
dependence with a floor at low wind,
\begin{equation}
  f_6 = \begin{cases} 0.25\,\nu_4, & I_8<0.25\\ \nu_4\,I_8, & I_8\ge0.25\end{cases}
  \qquad
  f_{6,\text{CO2}} = \begin{cases} 0.25\,\nu_{4\text{CO2}}, & I_8<0.25\\ \nu_{4\text{CO2}}\,I_8, & I_8\ge0.25,\end{cases}
\end{equation}
with separate heat ($\nu_4$) and \ce{CO2} ($\nu_{4\text{CO2}}$) leakage coefficients. The total
exchange rates entering the balances are
\begin{align}
  f_2 &= \eta_6\,f_5 \ \ \text{(if } \eta_7\ge\eta_8\text{) or } \eta_6 U_1 f_5, &
  f_3 &= \frac{U_7\,\phi_8}{\alpha_6}\ \text{(forced fans)}, &
  f_4 &= \eta_6\,f_7,
\end{align}
with $\eta_6$ the ventilation power-reduction factor and $U_7$ the forced-ventilation actuator.

\paragraph{Sensible-heat loss by ventilation.}
\begin{equation}
  h_7 \;=\; \rho_3\,\alpha_5\,\big(f_2 + f_4 + f_3 + \tfrac12 f_6\big)\,(\Tair - I_5).
\end{equation}
The roof term $f_4$ is included on the same footing as the side ($f_2$), forced ($f_3$) and
half-leakage ($\tfrac12 f_6$) terms: any air-exchange rate carries sensible heat regardless of what
drives it. (Omitting $f_4$ here makes a calm, sunny greenhouse unable to shed its solar load.)

\subsection{Sensible-heat and latent-heat fluxes of the air}\label{sec:heatfluxes}
\begin{align}
  h_1 &= 2\,\alpha_4\,I_1\,(\Tcan - \Tair) &&\text{(canopy}\to\text{air convection)},\\
  h_2 &= \frac{U_2\phi_7}{\alpha_6}\Big(\rho_3\alpha_5 I_5 - \gamma_2\rho_3\,\eta_5(\phi_5-\phi_6)\Big) &&\text{(fan--pad inlet)},\\
  h_4 &= n_{\text{pipes}}\,\big(1.99\,\pi\,\phi_1\,\gamma_1\,|I_3-\Tair|^{0.32}\big)(I_3-\Tair) &&\text{(heating pipe)},\\
  h_{11} &= \frac{2\,\nu_7\,(\Tair - I_7)}{\phi_2+\nu_8} &&\text{(soil-path conduction)},\\
  h_{12} &= (\eta_{15}+\eta_{16})\,\alpha_{12}\,U_{12} &&\text{(lamp reflector + direct heat)},\\
  l_2 &= \frac{\gamma_2\,U_9\,\phi_9}{\alpha_6} &&\text{(fog evaporation, latent)}.
\end{align}
$h_4$ is the natural-convection pipe--air exchange (a $|\Delta T|^{0.32}$ Nusselt correlation),
$U_2$ the pad actuator, $U_9$ the fog actuator.

\subsection{\ce{CO2} and vapour fluxes}\label{sec:co2vap}
\ce{CO2} fluxes (\si{mg.m^{-2}.s^{-1}}):
\begin{align}
  o_1 &= \frac{\eta_{13}\,U_4\,\lambda_4}{\alpha_6}, &
  o_2 &= \frac{U_{10}\,\psi_2}{\alpha_6}, &
  o_3 &= \frac{U_2\phi_7}{\alpha_6}(\Cair-I_{10}),\\
  o_4 &= A, &
  o_5 &= (\Cair-I_{10})\,\big(f_2+f_3+f_4+f_{6,\text{CO2}}\big), &&
\end{align}
with $I_{10}$ the outside \ce{CO2}, $U_4$ the air-heater actuator, $\lambda_4$ its capacity,
$\eta_{13}$ the \ce{CO2} produced per joule of combustion, $U_{10}$ the dosing valve and $\psi_2$
its capacity. Vapour fluxes (\si{Pa}-rate contributions, before $\kappa_3$):
\begin{align}
  p_1 &= q_1(q_2-\Vair), & p_2 &= \rho_3\frac{U_2\phi_7}{\alpha_6}\big(\eta_5(\phi_5-\phi_6)+\phi_6\big),\\
  p_3 &= \frac{U_9\phi_9}{\alpha_6}, & p_4 &= \frac{\eta_{12}\,U_4\,\lambda_4}{\alpha_6},\\
  p_5 &= \frac{\psi_1}{\omega_2}\Big(\frac{\Vair}{\Tair}-\frac{I_{11}}{I_5}\Big)\big(f_2+f_3+f_4+f_6\big),
      & p_6 &= \frac{U_2\phi_7}{\alpha_6}\frac{\psi_1}{\omega_2}\frac{\Vair}{\Tair},
\end{align}
with $I_{11}$ the outside vapour pressure. Condensation occurs only when the air is supersaturated
relative to the receiving surface (a one-sided switch):
\begin{align}
  p_7 &= \begin{cases}0,&\Vair<q_6\\ 6.4\!\times\!10^{-9}(\Vair-q_6),&\text{else}\end{cases} &&\text{(cooler)},\\
  p_8 &= \begin{cases}0,&\Vair<\mathrm{Pws}(T_{\text{scr}})\\ 6.4\!\times\!10^{-9}\!\cdot\!1.7\,|\Tair-\mathrm{Pws}(T_{\text{scr}})|^{0.33}(\Vair-\mathrm{Pws}(T_{\text{scr}})),&\text{else}\end{cases} &&\text{(screen)},\\
  p_9 &= \begin{cases}0,&\Vair<\mathrm{Pws}(\Tair)\\ 6.4\!\times\!10^{-9}\!\cdot\!1.2\!\cdot\!10^{3}\,|\Tcan-\Tair|^{0.33}(\Vair-\mathrm{Pws}(\Tair)),&\text{else}\end{cases} &&\text{(cover)},
\end{align}
where $T_{\text{scr}}=\tfrac12(I_5+\Tair)$ is the screen temperature.

\subsection{Floor / soil thermal mass}\label{sec:floor}
\paragraph{Principle.} The greenhouse floor and the top soil layer form a large \emph{thermal
buffer}: they absorb heat by day and release it at night, damping the air-temperature swing. We
model this with a single lumped floor temperature $\Tflo$ of areal heat capacity $C_{\text{soil}}$,
exchanging with the air above (conductance $k_{\text{ground}}$) and with a fixed deep-soil
temperature $T_{\text{deep}}$ below (same conductance):
\begin{align}
  Q_{af} &= k_{\text{ground}}(\Tair - \Tflo), &
  Q_{fg} &= k_{\text{ground}}(\Tflo - T_{\text{deep}}),\\
  \dot{T}_{2,\text{floor}} &= -\frac{Q_{af}}{\phi_2\rho_3\alpha_5}, &
  \dd{\Tflo}{t} &= \frac{Q_{af}-Q_{fg}}{C_{\text{soil}}}.
  \label{eq:floor}
\end{align}
The first equation returns the extra air-temperature rate added in Eq.~\eqref{eq:dT2}; the second is
the floor state equation. A single conductance $k_{\text{ground}}$ (rather than a separate
air--floor and floor--deep pair) is used because only the combination is identifiable from data.

\subsection{Psychrometrics and sky temperature}\label{sec:psychro}\label{sec:sky}
\paragraph{Saturation vapour pressure (Arden Buck).} For air/leaf temperature $T$ in kelvin, with
$t_c=T-273.15$ in \si{\degreeCelsius},
\begin{equation}
  \mathrm{Pws}(T)=\begin{cases}
    611.21\,\exp\!\Big[\big(18.678-\tfrac{t_c}{234.5}\big)\tfrac{t_c}{257.14+t_c}\Big], & t_c>0,\\[4pt]
    611.15\,\exp\!\Big[\big(23.036-\tfrac{t_c}{333.7}\big)\tfrac{t_c}{279.82+t_c}\Big], & t_c\le 0.
  \end{cases}
\end{equation}
Relative humidity and vapour-pressure deficit follow as
\begin{equation}
  \mathrm{RH}=100\,\frac{\Vair}{\mathrm{Pws}(\Tair)}\ [\%], \qquad
  \mathrm{VPD}=\mathrm{Pws}(\Tair)-\Vair\ [\si{Pa}].
\end{equation}
\paragraph{Effective sky temperature (Berdahl--Martin).} The radiative sky temperature $I_4$ is
colder than the air on clear nights and approaches the air temperature under overcast skies. From
air temperature $t_c$ [\si{\degreeCelsius}], relative humidity $\mathrm{RH}$ [\%] and total cloud
cover $c\in[0,1]$, the dew point (Magnus) and clear-sky emissivity are
\begin{align}
  \gamma_m &= \frac{a\,t_c}{b+t_c}+\ln\frac{\mathrm{RH}}{100}, \quad
  T_{\text{dew}} = \frac{b\,\gamma_m}{a-\gamma_m}, \quad (a=17.27,\ b=237.7),\\
  \epsilon_{\text{clear}} &= \mathrm{clip}\!\Big(0.711 + 0.56\tfrac{T_{\text{dew}}}{100} + 0.73\big(\tfrac{T_{\text{dew}}}{100}\big)^2,\ 0.6,\ 1\Big),\\
  \epsilon_{\text{sky}} &= \epsilon_{\text{clear}} + (1-\epsilon_{\text{clear}})\,c, \qquad
  I_4 = (t_c+273.15)\,\epsilon_{\text{sky}}^{1/4}.
\end{align}
This replaces a fixed cold sky and is the dominant control on night-time radiative loss.

%======================================================================
\section{Photosynthesis subsystem}\label{sec:photo}
%======================================================================

\paragraph{What this section contains.} The biochemical leaf photosynthesis model
(Farquhar--von~Caemmerer--Berry, FvCB), the temperature dependence of its parameters, and the
\textbf{canopy light-extinction model} that scales the leaf rate up to the whole canopy. The output
is the canopy gross assimilation rate $A$ that appears as $o_4$ in the \ce{CO2} balance
Eq.~\eqref{eq:dC1} and, integrated over the day (Eq.~\eqref{eq:Pg}), as the crop's carbon supply.
We place the light-extinction (Beer's-law canopy) model \emph{here}, because it is the bridge from
leaf biochemistry to the canopy source term; the same extinction physics also governs the PAR/NIR
absorption in the climate radiation submodel (Section~\ref{sec:rad}).

\subsection{The FvCB biochemical model}
\paragraph{Principle.} In C$_3$ plants (cucumber included) carbon fixation is limited by the slower
of three processes: (i) the maximum rate of the Rubisco enzyme, $W_c$; (ii) the rate at which the
light reactions regenerate RuBP, i.e.\ the electron-transport-limited rate $W_j$; and (iii) the rate
at which triose phosphates are used, $W_p$. The net assimilation is the minimum of the three,
corrected for photorespiration (through the \ce{CO2} compensation point $\Gamma^{\!*}$) and dark
respiration $R_d$. The inputs are leaf temperature $T$ (here the air temperature $\Tair$),
intercellular \ce{CO2} (taken proportional to the air \ce{CO2} $\Cair$) and absorbed irradiance
(from the canopy PAR $I_9$).

\paragraph{Inputs and conversions.} With incident radiation $I_w$ [\si{W.m^{-2}}] and air \ce{CO2}
$C$ [\si{mg.m^{-3}}],
\begin{equation}
  I = 4.6\times10^{-6}\,I_w \quad[\si{mol\,photons.m^{-2}.s^{-1}}], \qquad
  C_i = 0.509\times10^{-6}\,C \quad[\si{mol.m^{-3}}].
\end{equation}
\paragraph{\ce{CO2} compensation point and electron transport.}
\begin{align}
  \Gamma^{\!*} &= \frac{0.5\,O_a}{S_{c/o,25}\,\exp\!\big[\tfrac{T-298.15}{T}\cdot\tfrac{E_{S_{c/o}}}{298.15\,R}\big]},\\
  J &= \frac{(\alpha I + J_{\max}) - \sqrt{(\alpha I + J_{\max})^2 - 4\,\theta\,\alpha I\,J_{\max}}}{2\,\theta},
\end{align}
where $O_a$ is the \ce{O2} mole fraction, $S_{c/o,25}$ the Rubisco specificity and $E_{S_{c/o}}$ its
activation energy; $\alpha$ is the initial quantum-yield slope, $J_{\max}$ the maximum electron
transport rate and $\theta$ the curvature of the non-rectangular hyperbola relating $J$ to light.
\paragraph{Temperature-adjusted Rubisco kinetics.}
\begin{align}
  V_{c\max} &= V_{c\max,25}\,\frac{Q_{10,Vc}^{\,0.1(T-298.15)}}{1+\exp[0.128\,(T-315.15)]}, \\
  K_C &= K_{C,25}\,Q_{10,KC}^{\,0.1(T-298.15)}, \qquad
  K_O = K_{O,25}\,Q_{10,KO}^{\,0.1(T-298.15)},
\end{align}
with a peaked (deactivating) temperature response for $V_{c\max}$ and Michaelis constants $K_C,K_O$
for \ce{CO2} and \ce{O2}. The effective (stomatal+mesophyll) conductance is held at the constant
Bush (2023) value
\begin{equation}
  g_{\text{eff}} = \frac{0.3\cdot 0.25}{0.3+0.25} \approx 0.1364 \quad[\si{mol.m^{-2}.s^{-1}}].
\end{equation}
\paragraph{The three limiting rates.} Each limitation gives a quadratic in the assimilation rate
$A_r$; the physically admissible (lower) root is taken:
\begin{align}
  \text{Rubisco:}\quad & p_c = -\big(V_{c\max} + g_{\text{eff}}(C_i + K_C(1+O_a/K_O)) - R_d\big),\notag\\
  & q_c = g_{\text{eff}}\big(V_{c\max}\max(C_i-\Gamma^{\!*},0) - (C_i+K_C(1+O_a/K_O))R_d\big),\notag\\
  & A_{r,c} = \tfrac12\big(-p_c-\sqrt{p_c^2-4q_c}\big),\\[3pt]
  \text{Light:}\quad & p_j = -0.25\,J + R_d - g_{\text{eff}}(C_i+2\Gamma^{\!*}),\notag\\
  & q_j = 0.25\,g_{\text{eff}}\max(C_i-\Gamma^{\!*},0)\,J - g_{\text{eff}}(C_i+2\Gamma^{\!*})R_d,\notag\\
  & A_{r,j} = \tfrac12\big(-p_j-\sqrt{p_j^2-4q_j}\big),\\[3pt]
  \text{Product:}\quad & A_{r,p} = 1.5\,V_{c\max} - R_d.
\end{align}
The \textbf{leaf-level} net assimilation (per unit \emph{leaf} area) is $A_{\text{leaf}} =
\min(A_{r,j},A_{r,c},A_{r,p})$, and the \textbf{canopy} assimilation returned to the balances is
\begin{equation}
  A \;=\; 44010\;\cdot\;L_{\text{eff}}\;\cdot\;\min(A_{r,j},A_{r,c},A_{r,p}) \quad[\si{mg\,CO_2.m^{-2}.s^{-1}}],
  \label{eq:Acanopy}
\end{equation}
where the factor $44010$ converts \si{mol}$\to$\si{mg} of \ce{CO2}, and $L_{\text{eff}}$ is the
light-intercepting effective LAI defined next. For the cucumber crop under \ce{CO2} enrichment the
canopy operates predominantly in the \emph{light-limited} regime ($A_{r,j}$ smallest), so radiation
--- solar plus supplemental lamp PAR --- is the limiting resource.

\subsection{Canopy light-extinction model}\label{sec:lightext}
\paragraph{Principle (Beer--Lambert).} Leaves higher in the canopy shade those below, so irradiance
decays with cumulative leaf area $\ell$ (LAI measured from the top) as $I(\ell)=I_0\,e^{-k\ell}$,
with $k$ the canopy extinction coefficient ($k\approx0.7$ for cucumber). Two canopy-scaling schemes
are used.

\paragraph{(a) Big-leaf (effective-LAI) scaling.} The leaf rate is evaluated once at the
top-of-canopy light $I_0$ and multiplied by the \emph{light-intercepting effective LAI}
\begin{equation}
  L_{\text{eff}}(\mathrm{LAI}) \;=\; \frac{1-e^{-k\,\mathrm{LAI}}}{k},
  \label{eq:laieff}
\end{equation}
which is the vertical integral of the interception profile and \emph{saturates} at $1/k$ as
LAI$\to\infty$ (extra leaf intercepts almost no additional light). This $L_{\text{eff}}$ is the
factor used in Eq.~\eqref{eq:Acanopy} and, with the corresponding $\beta_1,\beta_2,\beta_3$
coefficients, underlies the PAR/NIR absorption of Section~\ref{sec:rad}.

\paragraph{(b) Layered (vertical-profile) scaling.} A more faithful scheme integrates the
\emph{leaf} rate down the light gradient. Discretising the canopy into leaf layers of area $a_i$ at
cumulative depth $\ell_i$ (mid-layer),
\begin{equation}
  A \;=\; 44010\;\sum_i a_i\; A_{\text{leaf}}\!\big(I_0\,e^{-k\,\ell_i}\big).
  \label{eq:layered}
\end{equation}
When the leaf response is linear in light (the light-limited regime) Eq.~\eqref{eq:layered} reduces
exactly to the big-leaf form~\eqref{eq:Acanopy}--\eqref{eq:laieff}; when upper leaves are
light-saturated it correctly predicts a \emph{higher} canopy assimilation, because shaded lower
leaves operate at their more efficient (steeper) part of the light-response curve. Scheme (b) is the
one used by the leaf-cohort crop model (Section~\ref{sec:cohort}), where the layers are the leaf-age
cohorts; scheme (a) is used by the single-pool crop model (Section~\ref{sec:biggrow}).

%======================================================================
\section{Crop-growth subsystem}\label{sec:growth}
%======================================================================

\paragraph{What this section contains.} The daily crop model: how the day's assimilate $P_g$ is
partitioned into maintenance, organ growth and fruit; how development proceeds in thermal time; the
\textbf{flowering / fruit-set equations}; fruit growth and harvest; and the leaf-area update. Two
variants share the same physiology and differ only in how the canopy is represented: a
\emph{single-pool} model (Section~\ref{sec:biggrow}) and a \emph{leaf-age-cohort / vertical-canopy}
model (Section~\ref{sec:cohort}). The \textbf{planting density} and its role are defined in
Section~\ref{sec:density}. The governing biological principle is \textbf{source--sink carbon
partitioning}: a common assimilate pool (the source) is distributed among competing organs (the
sinks) in proportion to their potential growth (sink strength), limited by supply when the source is
scarce.

\subsection{State variables and thermal time}\label{sec:cropstate}
The single-pool crop state is
\begin{equation}
  \big(W_{\text{leaf}},\ W_{\text{stem}},\ W_{\text{root}},\ \mathrm{LAI},\ \{\text{fruit cohorts}\},\ Y_{DW},\ Y_{FW},\ n_{\text{node}}\big),
\end{equation}
with organ dry weights $W_\bullet$ [\si{g.m^{-2}}], leaf area index LAI, a list of fruit cohorts, the
cumulative harvested fruit dry ($Y_{DW}$, \si{g.m^{-2}}) and fresh ($Y_{FW}$, \si{kg.m^{-2}}) weight,
and the main-stem node number $n_{\text{node}}$. Each \textbf{fruit cohort} carries a development
stage $\mathrm{dev}\in[0,1]$ (its thermal age), a dry weight $W$ [\si{g.m^{-2}}] and a number of
fruits $n$ [\si{fruits.m^{-2}}].

\paragraph{Thermal time.} Development is driven by temperature above a base $T_{\text{base}}$
(\textbf{degree-days}): with daily-mean air temperature $\overline{\Tair}$ [\si{\degreeCelsius}],
\begin{equation}
  \Delta T \;=\; \max\!\big(\overline{\Tair}-T_{\text{base}},\,0\big).
\end{equation}
Nodes appear on the growing head at rate $r_{\text{node}}$ until \emph{topping} on day
$d_{\text{top}}$, and each fruit cohort ages by $\Delta\mathrm{dev}=\Delta T/\mathrm{DD}_{\text{fruit}}$
per day, where $\mathrm{DD}_{\text{fruit}}$ is the degree-days from set to ripe:
\begin{equation}
  n_{\text{node}} \mathrel{+}= r_{\text{node}}\,\Delta T \ \ (\text{if } d<d_{\text{top}}),\qquad
  \mathrm{dev} \mathrel{+}= \frac{\Delta T}{\mathrm{DD}_{\text{fruit}}}\ \ \text{(all cohorts)}.
\end{equation}

\subsection{Respiration, supply and growth conversion}\label{sec:resp}
\paragraph{Maintenance respiration.} Living tissue consumes carbohydrate to maintain itself at a
rate that doubles roughly every \SI{10}{\degreeCelsius} (a $Q_{10}$ law):
\begin{align}
  q_{10} &= Q_{10,\text{resp}}^{\,(\overline{\Tair}-25)/10},\\
  R_m &= q_{10}\big(k_{m,\text{leaf}}W_{\text{leaf}} + k_{m,\text{stem}}W_{\text{stem}} + k_{m,\text{root}}W_{\text{root}} + k_{m,\text{fruit}}W_{\text{fruit}}\big),
\end{align}
with organ-specific maintenance coefficients $k_{m,\bullet}$ [\si{g\,CH_2O.g^{-1}.d^{-1}}] and
$W_{\text{fruit}}=\sum_{\text{cohorts}}W$. \paragraph{Assimilate available for growth.} What remains
of the day's gross assimilate after maintenance is converted to structural dry matter with a growth
(conversion) efficiency $1/\mathrm{asrq}$:
\begin{equation}
  C_{\text{av}} = \max(P_g - R_m,\,0), \qquad
  \Delta \mathrm{DM}_{\text{pot}} = \frac{C_{\text{av}}}{\mathrm{asrq}} \quad[\si{g\,DM.m^{-2}.d^{-1}}],
  \label{eq:supply}
\end{equation}
where $\mathrm{asrq}$ (assimilate required per unit dry matter, $>1$) accounts for growth
respiration. $\Delta\mathrm{DM}_{\text{pot}}$ is the potential (source-side) dry-matter increment.

\subsection{Sinks, source/sink limitation and allocation}\label{sec:alloc}
\paragraph{Sink strengths.} The vegetative sink fills the canopy toward a density-dependent cap
$L_{\text{cap}}$ and switches off as the cap is approached; each fruit cohort's sink follows a
bell-shaped growth curve. With the density factor $\mathrm{d}=\mathrm{stem\_density}/\mathrm{d}_{\text{ref}}$
(Section~\ref{sec:density}),
\begin{align}
  L_{\text{cap}} &= \mathrm{LAI}_{\max}\,\mathrm{d}, &
  S_{\text{veg}} &= \max\!\Big(\mathrm{veg\_sink}_{\max}\,\mathrm{d}\,\big(1-\tfrac{\mathrm{LAI}}{L_{\text{cap}}}\big),\,0\Big),\\
  S_{\text{fruit},i} &= n_i\,W_{f\max}\,\varsigma(\mathrm{dev}_i)\,\Delta\mathrm{dev}, &
  \varsigma(\mathrm{dev}) &= 6\,\mathrm{dev}\,(1-\mathrm{dev}),
\end{align}
where $W_{f\max}$ is the potential mature fruit dry weight and the \emph{sink-shape}
$\varsigma(\mathrm{dev})=6\,\mathrm{dev}(1-\mathrm{dev})$ (the derivative of the smoothstep
$\mathrm{dev}^2(3-2\mathrm{dev})$) integrates to unity over $\mathrm{dev}\in[0,1]$, so a cohort's
lifetime potential growth totals $n\,W_{f\max}$. \paragraph{Allocation.} The total demand is
$S_{\text{tot}}=S_{\text{veg}}+\sum_i S_{\text{fruit},i}$. The realised increment is
\emph{source-limited} when supply is scarce and \emph{sink-limited} when supply is ample:
\begin{equation}
  \Delta\mathrm{DM} = \min\big(\Delta\mathrm{DM}_{\text{pot}},\,S_{\text{tot}}\big),
\end{equation}
and is split among organs in proportion to their sink share:
\begin{align}
  \Delta\mathrm{DM}_{\text{veg}} &= \Delta\mathrm{DM}\,\frac{S_{\text{veg}}}{S_{\text{tot}}}, &
  W_{\text{leaf}} &\mathrel{+}= f_{\text{leaf}}\,\Delta\mathrm{DM}_{\text{veg}},\\
  W_{\text{stem}} &\mathrel{+}= f_{\text{stem}}\,\Delta\mathrm{DM}_{\text{veg}}, &
  W_{\text{root}} &\mathrel{+}= f_{\text{root}}\,\Delta\mathrm{DM}_{\text{veg}},\\
  W_i &\mathrel{+}= \Delta\mathrm{DM}\,\frac{S_{\text{fruit},i}}{S_{\text{tot}}} \quad\text{(each fruit cohort)}, &&
\end{align}
with vegetative partitioning fractions $f_{\text{leaf}}+f_{\text{stem}}+f_{\text{root}}=1$.

\subsection{Flowering: carbohydrate-regulated fruit set}\label{sec:flowering}
\paragraph{Principle.} New fruits are set continuously once the plant is old enough
($d\ge d_{\text{set,start}}$), but the \emph{rate} of setting is regulated by the plant's
carbohydrate status: when the standing fruit load already demands more carbon than the source can
supply, further set is suppressed (young fruit abort / fail to set); once those fruits are harvested
and the load drops, set resumes. Combined with the $\sim$2-week fruit-growth delay
($\mathrm{DD}_{\text{fruit}}$), this supply/demand feedback produces the characteristic weekly
harvest \emph{flushes}. Define the instantaneous supply/demand ratio
\begin{equation}
  s \;=\; \frac{\Delta\mathrm{DM}_{\text{pot}}}{S_{\text{tot,pre}} + 10^{-9}},
  \qquad S_{\text{tot,pre}} = S_{\text{veg}} + \textstyle\sum_i S_{\text{fruit},i},
\end{equation}
evaluated \emph{before} the new cohort is added. The set fraction is a clipped linear ramp between a
lower and upper carbohydrate threshold,
\begin{equation}
  f_{\text{set}} \;=\; \mathrm{clip}\!\Big(\frac{s - s_{\text{low}}}{s_{\text{high}} - s_{\text{low}}},\,0,\,1\Big),
  \label{eq:setfrac}
\end{equation}
so that below $s_{\text{low}}$ no fruit sets and above $s_{\text{high}}$ set is maximal. On each day
with $d_{\text{set,start}}\le d<d_{\text{top}}$ a new fruit cohort is created with number
\begin{equation}
  n_{\text{set}} \;=\; \mathrm{set\_rate}\,\cdot\,\mathrm{d}\,\cdot\, f_{\text{set}}
  \qquad(\text{single-pool model}),
  \label{eq:nset}
\end{equation}
where $\mathrm{set\_rate}$ is the potential fruits per unit area per day and $\mathrm{d}$ the density
factor. (In the cohort model, Section~\ref{sec:cohort}, set is driven by node appearance rather than
calendar day; see Eq.~\eqref{eq:nset_cohort}.)

\subsection{Fruit growth, harvest and leaf-area update}\label{sec:harvest}
A cohort is harvested when it reaches $\mathrm{dev}\ge1$; its dry weight is banked and converted to
fresh weight through the fruit dry-matter content $\mathrm{DMC}_{\text{fruit}}$:
\begin{equation}
  Y_{DW}\mathrel{+}=W, \qquad
  Y_{FW}\mathrel{+}=\frac{W/\mathrm{DMC}_{\text{fruit}}}{1000}\ \ [\si{kg\,FW.m^{-2}}].
\end{equation}
Finally the leaf area is refreshed from the leaf dry weight through the specific leaf area (SLA,
leaf area per unit leaf dry weight) and capped (see pruning, Section~\ref{sec:pruning}):
\begin{equation}
  \mathrm{LAI} = \mathrm{SLA}\cdot W_{\text{leaf}}.
  \label{eq:lai_sla}
\end{equation}

\subsection{Single-pool canopy model}\label{sec:biggrow}
The model of Sections~\ref{sec:cropstate}--\ref{sec:harvest} \emph{is} the single-pool model: leaves
form one lumped pool $W_{\text{leaf}}$, the canopy source uses the big-leaf light scaling
(Eq.~\eqref{eq:laieff}), and pruning is a cap on total leaf area (Section~\ref{sec:pruning}). It is
compact and fast, but cannot represent the vertical age/light structure of a real high-wire cucumber
canopy --- which motivates the cohort model.

\subsection{Leaf-age-cohort (vertical-canopy) model}\label{sec:cohort}
\paragraph{Principle.} A high-wire cucumber produces new leaves at the head while old leaves at the
base become deeply shaded, senesce, and are periodically removed by the grower. The cohort model
represents the canopy as a stack of \textbf{leaf-age cohorts}, each with an area and a thermal age,
so that (i) photosynthesis is integrated layer-by-layer down the light gradient
(Eq.~\eqref{eq:layered}), (ii) each cohort's maintenance respiration is charged on its own mass, and
(iii) de-leafing removes the oldest, most-shaded cohorts. This yields a physically meaningful
optimal LAI and a faithful account of de-leafing that the single pool cannot provide.

\paragraph{Leaf production (node-driven).} Rather than filling a cap, leaf area is produced in step
with node appearance. With leaf area demanded per node $\ell_{\text{node}}$ and $\Delta n_{\text{node}}
= r_{\text{node}}\Delta T$ the day's new nodes,
\begin{equation}
  \text{leaf-area demand} = \ell_{\text{node}}\,\mathrm{d}\,\Delta n_{\text{node}},\qquad
  S_{\text{veg}} = \frac{\ell_{\text{node}}\,\mathrm{d}\,\Delta n_{\text{node}}/\mathrm{SLA}}{f_{\text{leaf}}},
\end{equation}
the new top cohort receiving $f_{\text{leaf}}\,\Delta\mathrm{DM}_{\text{veg}}\,\mathrm{SLA}$ of area.
\paragraph{Node-driven flowering.} Fruit set is triggered per node flush rather than per day, with the
same carbohydrate regulation $f_{\text{set}}$ of Eq.~\eqref{eq:setfrac}:
\begin{equation}
  n_{\text{set}} = \mathrm{set\_rate}\,\cdot\,\mathrm{d}\,\cdot\,\Delta n_{\text{node}}\,\cdot\, f_{\text{set}}
  \qquad(\text{cohort model}).
  \label{eq:nset_cohort}
\end{equation}
\paragraph{Senescence and de-leafing.} A cohort is dropped by natural senescence once its thermal age
exceeds the leaf lifespan $\mathrm{life}_{\text{leaf}}$ [degree-days]. Management de-leafing removes
the \emph{oldest} cohorts down to a working target LAI $L_{\text{deleaf}}$: with
$\mathrm{LAI}=\sum_i a_i$, cohorts are trimmed oldest-first until $\mathrm{LAI}\le L_{\text{deleaf}}$.
Because the removed leaves are the deepest and most shaded, they contribute almost no assimilation
(Eq.~\eqref{eq:layered}) yet still respired, so their removal is close to cost-free in carbon and
sheds their maintenance load.

\subsection{Planting density}\label{sec:density}
\paragraph{Where it lives and what it does.} Planting density (stems per \si{m^2}) is a
\emph{crop-structural / management} parameter, so it is defined here, in the growth subsystem; it is
also a control decision, discussed in Section~\ref{sec:cropmgmt}. All the per-area capacities of the
growth model (the vegetative sink $\mathrm{veg\_sink}_{\max}$, the fruit-set rate $\mathrm{set\_rate}$,
the leaf-area cap $\mathrm{LAI}_{\max}$ and the per-node leaf demand $\ell_{\text{node}}$) are defined
at a \emph{reference} density $\mathrm{d}_{\text{ref}}$ and scale linearly with the actual density
through the factor
\begin{equation}
  \mathrm{d} \;=\; \frac{\mathrm{stem\_density}}{\mathrm{d}_{\text{ref}}},
\end{equation}
which appears in the sink and set equations above. At $\mathrm{stem\_density}=\mathrm{d}_{\text{ref}}$
the factor is unity, so all reference calibrations are recovered exactly. A denser stand closes the
canopy (raises LAI) faster and intercepts light earlier; because interception saturates
(Eq.~\eqref{eq:laieff}) while maintenance respiration keeps growing with leaf mass, there is a
finite \emph{optimal} density, and its value is governed primarily by the light-extinction
coefficient $k$ (Section~\ref{sec:lightext}). The initial canopy is also scaled by density,
$\mathrm{LAI}_0\!\leftarrow\!\mathrm{LAI}_0\,\mathrm{d}$.

%======================================================================
\section{Control variables and management strategies}\label{sec:control}
%======================================================================

\paragraph{What this section contains.} The actuators that drive the model, the setpoint
programme that a grower (or an automatic controller) specifies, the low-level feedback controllers
that translate setpoints into actuator positions, and the discrete crop-management strategies
(pruning/de-leafing, topping, density scheduling). The controlled variable is the greenhouse
\emph{air} temperature $\Tair$ (the physically measured quantity), not the canopy temperature.

\subsection{Actuators}
Twelve actuators $U_1,\dots,U_{12}$ enter the climate right-hand side (Table~\ref{tab:controls}).
In the closed-loop twin only a subset is active: the thermal/energy screen $U_1$, the side vent
$U_6$, the roof vent $U_8$, the fog $U_9$, the \ce{CO2} valve $U_{10}$ and the lamps $U_{12}$; the
heating pipe enters as its temperature $I_3$. The pad/heater/forced-ventilation actuators
($U_2,U_3,U_4,U_5,U_7$) are available but set to zero in the reference configuration.

\subsection{Setpoint programme (the action surface)}
A control programme is a set of within-day setpoint schedules:
\begin{itemize}[leftmargin=1.6em,itemsep=1pt]
  \item day and night \textbf{air-temperature} setpoints $T_{\text{set}}$ (a day/night step schedule);
  \item \textbf{\ce{CO2}} setpoint $C_{\text{set}}$;
  \item \textbf{lamp} on/off schedule (start hour, duration $\Rightarrow$ end hour);
  \item ventilation \textbf{proportional band} $p_{\text{band}}$ and dead-band \emph{offset} above
  the temperature setpoint;
  \item the outside-temperature threshold $T_{\text{out,max}}$ (screen-closing) and radiation
  threshold $\mathrm{Rad}_{\text{scr}}$;
  \item the \textbf{VPD floor} $\mathrm{VPD}_{\min}$ for humidity control.
\end{itemize}

\subsection{Heating controller (PI $\to$ pipe temperature)}
A leaky-integral proportional--integral controller on the temperature error
$e_T=T_{\text{set}}-\Tair$ sets a heating fraction $U_{\text{heat}}\in[0,1]$, which blends the
maximum pipe temperature $T_{\text{pipe}}^{\max}$ with the air temperature to give the pipe
temperature $I_3$:
\begin{equation}
  U_{\text{heat}} = \mathrm{clip}\big(K_p^{T} e_T + K_i^{T} e_T^{\text{int}},\,0,\,1\big), \qquad
  I_3 = U_{\text{heat}}\,T_{\text{pipe}}^{\max} + (1-U_{\text{heat}})\,\Tair,
\end{equation}
where the integral state obeys $\dot{e}_T^{\text{int}} = e_T - \lambda_T e_T^{\text{int}}$ (a
\emph{leaky} integrator with decay $\lambda_T$), \emph{frozen} under actuator saturation
(anti-windup, Section~\ref{sec:windup}). $I_3$ then drives the pipe heat flux $h_4$
(Section~\ref{sec:heatfluxes}) and the pipe FIR $r_6$.

\subsection{\ce{CO2} dosing controller (PI $\to$ valve)}
Analogously, on the error $e_C=C_{\text{set}}-\Cair$,
\begin{equation}
  U_{10} = \mathrm{clip}\big(K_p^{C} e_C + K_i^{C} e_C^{\text{int}},\,0,\,1\big),
\end{equation}
opening the external \ce{CO2} valve (source $o_2=U_{10}\psi_2/\alpha_6$), with the same leaky-integral
and anti-windup treatment.

\subsection{Ventilation controller (proportional band $\to$ side/roof)}
Ventilation opens on a proportional band above the temperature setpoint plus a dead-band offset:
\begin{equation}
  U_{\text{vent}}=\begin{cases}
    0, & \Tair < T_{\text{set}}+\text{ofs},\\[2pt]
    \dfrac{\Tair - T_{\text{set}} - \text{ofs}}{p_{\text{band}}}, & T_{\text{set}}+\text{ofs}\le \Tair < T_{\text{set}}+\text{ofs}+p_{\text{band}},\\[6pt]
    1, & \text{otherwise},
  \end{cases}
\end{equation}
split so the side vent opens first and the roof vent follows:
\begin{equation}
  U_6 = 2\min(U_{\text{vent}},0.5), \qquad U_8 = 2\max(0,\,U_{\text{vent}}-0.5).
\end{equation}

\subsection{Thermal / energy-screen controller}
The screen is deployed to retain heat when it is dark (radiation $I_2<\mathrm{Rad}_{\text{scr}}$) or
cold outside ($I_5<T_{\text{out,max}}$), and is cracked open to shed moisture when the air is too
humid (VPD below the floor):
\begin{equation}
  U_1 = \begin{cases}
    0.85, & \text{closed and } \mathrm{VPD}<\mathrm{VPD}_{\min}\ \text{(crack to dehumidify)},\\
    1, & \text{closed (dark or cold) and not too humid},\\
    0, & \text{open (day)}.
  \end{cases}
\end{equation}
The screen's FIR emissivity is the low-emissivity \emph{energy-screen} value $\epsilon_{scr}$
(Section~\ref{sec:rad}), distinct from a full-emissivity shade screen, so that deploying it genuinely
retains long-wave heat.

\subsection{Humidity control by vapour-pressure deficit}
Humidity is regulated on the \textbf{VPD} (Section~\ref{sec:psychro}), not relative humidity, because
VPD is the physiologically and physically meaningful driver of transpiration and condensation and is
temperature-independent. When $\mathrm{VPD}=(\mathrm{Pws}(\Tair)-\Vair)/1000$ [\si{kPa}] falls below
the floor $\mathrm{VPD}_{\min}$ (air too close to saturation, condensation/disease risk), a
proportional dehumidification demand is added to the vents (and the screen is cracked as above):
\begin{equation}
  U_{\text{hv}} = \mathrm{clip}\!\Big(\frac{\mathrm{VPD}_{\min}-\mathrm{VPD}}{\text{band}},\,0,\,1\Big)\,v_{\max},
  \qquad U_6\leftarrow\max(U_6,U_{\text{hv}}),\ \ U_8\leftarrow\max(U_8,U_{\text{hv}}).
\end{equation}

\subsection{Supplemental lighting}
The lamps contribute PAR to photosynthesis (added to the canopy light $I_w$ in the layered source,
Section~\ref{sec:lightext}), NIR heat to the canopy ($r_5$, via $\eta_{14}$), and reflector/direct
heat to the air ($h_{12}$, via $\eta_{15},\eta_{16}$), at an electrical intensity $\alpha_{12}$ per
unit area when on ($U_{12}=1$). Because lamp electricity is the dominant operating cost and its value
depends on the free solar light available and on time-of-use tariffs, the lamp \emph{schedule}
(start hour and duration) is a primary management lever.

\subsection{Crop-management strategies}\label{sec:cropmgmt}
\paragraph{Pruning / de-leafing.}\label{sec:pruning} Leaf area is limited to the smaller of the
density cap $L_{\text{cap}}$ and a management de-leaf target $L_{\text{deleaf}}$. In the single-pool
model this caps $W_{\text{leaf}}$ (Eq.~\eqref{eq:lai_sla}); in the cohort model it removes the oldest,
most-shaded cohorts down to $L_{\text{deleaf}}$ (Section~\ref{sec:cohort}). De-leafing sheds the
maintenance-respiration burden of leaves that intercept little light, without touching stems or fruit.
\paragraph{Topping.} Node appearance (vegetative growth of the head) stops at the topping day
$d_{\text{top}}$, after which the plant only ripens its existing fruit load.
\paragraph{Density scheduling (interplanting / plant separation).} The planting density
(Section~\ref{sec:density}) may be changed during the season: planted densely for fast early canopy
closure, then thinned (stems pulled) later. A density reduction on a given day scales the standing
leaf area and the standing (unharvested) fruit load by the ratio of new-to-old density (the pulled
stems leave with their leaves and green fruit), and all subsequent per-area production runs at the
reduced density.

\subsection{Anti-windup}\label{sec:windup}
Because the PI integrators are ODE states, they can \emph{wind up} while an actuator is saturated. The
integral is therefore frozen (its accumulation term set to the leak only) whenever the actuator is at
a bound \emph{and} the error would push it further into saturation:
\begin{equation}
  \dot{e}^{\text{int}} = \begin{cases}
    -\lambda\,e^{\text{int}}, & \text{saturated in the driving direction},\\
    e - \lambda\,e^{\text{int}}, & \text{otherwise},
  \end{cases}
\end{equation}
for both the temperature ($e_T,\lambda_T$) and \ce{CO2} ($e_C,\lambda_C$) loops. This prevents the
runaway that windup causes once the leakage $\nu_4,\nu_{4\text{CO2}}$ is set to its physical value.
\appendix


\section{Climate-model constants}\label{app:climate}
Physical parameters of the climate subsystem (Section~\ref{sec:climate}). Values as configured in the reference twin; symbols follow the Vanthoor convention ($\alpha_i,\beta_i,\dots$).
\begin{longtable}{@{}p{0.16\textwidth} p{0.16\textwidth} p{0.20\textwidth} p{0.42\textwidth}@{}}
\toprule
\textbf{Name} & \textbf{Value} & \textbf{Units} & \textbf{Description}\\
\midrule
\endfirsthead
\toprule
\textbf{Name} & \textbf{Value} & \textbf{Units} & \textbf{Description}\\
\midrule
\endhead
\texttt{alpha1} & 1500 & J * K\textasciicircum\{\}-1 * m\textasciicircum\{\}-2 & Heat capacity of one square meter of the canopy\\
\texttt{alpha12} & 74 & W * m\textasciicircum\{\}-2 & Total lamp radiation per square meter \\
\texttt{alpha2} & 0.35 &  & Global NIR absorption coefficient of the canopy\\
\texttt{alpha3} & 0.3 & m\textasciicircum\{\}2*m\textasciicircum\{\}-2 & Surface of the heating pipe\\
\texttt{alpha4} & 5 & W * m\textasciicircum\{\}-2 * K\textasciicircum\{\}-1 & Convection heat exchange coefficient of canopy leaf to greenhouse air\\
\texttt{alpha5} & 1500 & J * K\textasciicircum\{\}-1 * kg\textasciicircum\{\}-1 & Specific heat capacity of greenhouse air\\
\texttt{alpha6} & 10000 & m\textasciicircum\{\}2 & Greenhouse floor surface area\\
\texttt{alpha7} & 0.5 &  & Global NIR absorption coefficient of the floor\\
\texttt{alpha8} & 1 &  & PAR absorption coefficient of the cover\\
\texttt{alpha9} & 1 &  & NIR absorption coefficient of the cover\\
\texttt{beta1} & 0.7 &  & Canopy extinction coefficient for PAR radiation\\
\texttt{beta2} & 0.7 &  & Extinction coefficient for PAR radiation reflected from the floor to the canopy\\
\texttt{beta3} & 0.27 &  & Canopy extinction coefficient for NIR radiation\\
\texttt{C\_soil} & 150000 & J m\textasciicircum\{\}-2 K\textasciicircum\{\}-1 & Heat capacity per floor area of the active soil/floor layer (thermal mass)\\
\texttt{delta1} & 5 & W * m\textasciicircum\{\}-2 & Radiation above the canopy that defines sunrise and sunset\\
\texttt{delta2} & 4.3 & W * m\textasciicircum\{\}-2 & Empirically determined parameter\\
\texttt{delta3} & 0.54 & W * m\textasciicircum\{\}-2 & Empirically determined parameter\\
\texttt{delta4} & 6.1e-07 & ppm\textasciicircum\{\}-2 & Coefficient of the CO2 transpiration in the day\\
\texttt{delta5} & 1.1e-11 & ppm\textasciicircum\{\}-2 & Coefficient of the CO2 transpiration in the night\\
\texttt{delta6} & 4.3e-06 & Pa\textasciicircum\{\}-2 & Coefficient of the vapour pressure in the day\\
\texttt{delta7} & 5.2e-06 & Pa\textasciicircum\{\}-2 & Coefficient of the vapour pressure in the night\\
\texttt{eps\_screen} & 0.4 &  & Energy-screen FIR emissivity to sky (aluminized/low-e). Separate from epsil5 so it only affects the screen-sky term r12 when the screen is deployed (U1>0); U1=0 on the venting calibration day, so calibration-safe.\\
\texttt{epsil1} & 0.88 &  & FIR emission coefficient of the heating pipe\\
\texttt{epsil2} & 1 &  & Canopy FIR emission coefficient\\
\texttt{epsil3} & 1 &  & Sky FIR emission coefficient\\
\texttt{epsil4} & 1 &  & Floor FIR emission coefficient\\
\texttt{epsil5} & 1 &  & Thermal screen FIR emission coefficient\\
\texttt{epsil6} & 0.44 &  & External cover FIR emission coefficient\\
\texttt{eta1} & 0.1 &  & Proportion of global radiation that is absorbed by greenhouse building elements\\
\texttt{eta10} & 0 &  & Shadow effect on the discharge coefficient\\
\texttt{eta11} & 0 &  & Effect of shadow on the global wind pressure coefficient\\
\texttt{eta12} & 4.43e-08 & kg\_vapour * J\textasciicircum\{\}-1 & Amount of vapor that is released when a joule of sensible energy is produced by the direct air heater\\
\texttt{eta13} & 0.057 & mg\_CO2 * J\textasciicircum\{\}-1 & Amount of CO2 that is released when a joule of sensible energy is produced by the direct air heater\\
\texttt{eta14} & 0.18 &  & Percentage lamps radiation that is NIR\\
\texttt{eta15} & 0.0074 &  & Porcentaje de la radiacion de las lamparas que es calor en el reflector\\
\texttt{eta16} & 0.39 &  & Porcentaje de la radiacion de las lamparas que es calor directo al aire del invernadero\\
\texttt{eta17} & 0.3626 &  & Percentage lamps radiation that is PAR\\
\texttt{eta2} & 0.5 &  & Ratio between PAR radiation and external global radiation\\
\texttt{eta3} & 0.5 &  & Ratio between NIR radiation and global external radiation\\
\texttt{eta4} & 0.554 & ppm * mg\textasciicircum\{\}-1 * m\textasciicircum\{\}3 & Conversion factor for CO2 of mg*m**-3 to ppm\\
\texttt{eta5} & 0.5 &  & Fan-pad system efficiency\\
\texttt{eta6} & 0.551 &  & Ventilation power reduction factor\\
\texttt{eta7} & 0.8 &  & Ratio between ceiling ventilation area and total ventilation area\\
\texttt{eta8} & 0.9 &  & Ratio between ceiling and total ventilation area, if there is no chimney effect\\
\texttt{gamma} & 65.8 & Pa * K\textasciicircum\{\}-1 & Psychometric constan\\
\texttt{gamma1} & 1.25 & m * m\textasciicircum\{\}-2 & Length of the heating pipe\\
\texttt{gamma2} & 2.45e+06 & J * kg\_water\textasciicircum\{\}-1 & Latent heat of water evaporation\\
\texttt{gamma3} & 275 & s * m\textasciicircum\{\}-1 & Strength of boundary layer of canopy for vapor transport\\
\texttt{gamma4} & 82 & s * m\textasciicircum\{\}-1 & Minimum stomatal resistance of the canopy\\
\texttt{gamma5} & -1 & W * m\textasciicircum\{\}-2 & Inverse of slope of the differentiable switch for the stomatal resistance model\\
\texttt{h\_soil} & 4 & W m\textasciicircum\{\}-2 K\textasciicircum\{\}-1 & Effective air<->floor convective heat-exchange coefficient (thermal-mass buffer)\\
\texttt{HEAT\_PIPE} & 363.15 &  & Temperatura del tubo de calentamiento\\
\texttt{k\_deep} & 2 & W m\textasciicircum\{\}-2 K\textasciicircum\{\}-1 & Conductance from the active floor layer to the fixed deep-soil temperature\\
\texttt{k\_ground} & 4.33 & W m\textasciicircum\{\}-2 K\textasciicircum\{\}-1 & Effective air<->ground conductance via the floor thermal-mass state (single, identifiable)\\
\texttt{lamb1} & 0 &  & Performance coefficient of the mechanical cooling system\\
\texttt{lamb2} & 0 & W & Electrical capacity of the mechanical cooling system\\
\texttt{lamb3} & 1 & W * m\textasciicircum\{\}-2 * K\textasciicircum\{\}-1 & Convective heat exchange coefficient between soil and greenhouse air\\
\texttt{lamb4} & 500000 & W & Heat capacity of direct air heater\\
\texttt{lamb5} & 18000 & m\textasciicircum\{\}2 & Cover surface\\
\texttt{lamb6} & 2.8 & W * m\_cover\textasciicircum\{\}-2 * K\textasciicircum\{\}-1 & Variable of heat exchange by convection between the roof and the outside air\\
\texttt{lamb7} & 1 & J * m\textasciicircum\{\}-3 * K\textasciicircum\{\}-1 & Variable of heat exchange by convection between the roof and the outside air\\
\texttt{lamb8} & 2 &  & Variable of heat exchange by convection between the roof and the outside air\\
\texttt{lamb9} & 0.05 &  & eficiencia del intercambio de calor entre las capas de aire\\
\texttt{n\_pipes} & 2 &  & Numero de tuberias de calentamiento\\
\texttt{nu1} & 1 &  & Shadowless discharge coefficient\\
\texttt{nu2} & 0.1 &  & Global wind pressure coefficient without shadow\\
\texttt{nu3} & 2352.9 & m\textasciicircum\{\}2 & Side surface of the greenhouse\\
\texttt{nu4} & 0.0001 &  & Heat leakage coefficient\\
\texttt{nu4CO2} & 0.0001 &  & CO2 leakage coefficient\\
\texttt{nu5} & 2000 & m\textasciicircum\{\}2 & Maximum ceiling ventilation area\\
\texttt{nu6} & 1 & m & Vertical dimension of a single open respirator\\
\texttt{nu7} & 0.85 & W * m\textasciicircum\{\}-1 * K\textasciicircum\{\}-1 & Soil thermal conductivity\\
\texttt{nu8} & 0.64 & m & Floor to ground distance\\
\texttt{omega1} & 9.81 & m * s\textasciicircum\{\}-2 & Gravity acceleration constant\\
\texttt{omega2} & 8314 & J * kmol\textasciicircum\{\}-1 * K\textasciicircum\{\}-1 & Molar gas constant\\
\texttt{omega3} & 0.0075 &  & Percentage of CO2 absorbed by the canopy\\
\texttt{phi1} & 0.051 & m & External diameter of the heating pipe\\
\texttt{phi2} & 4 & m & Average height of greenhouse air\\
\texttt{phi3} & 28.9647 & kg * kmol\textasciicircum\{\}-1 & Masa molar del aire\\
\texttt{phi4} & 0 & metros & Altitud del invernadero\\
\texttt{phi5} & 0.014 & kg\_water * kg\_air\textasciicircum\{\}-1 & Water vapor contained in the fan-pad system\\
\texttt{phi6} & 0.0079 & kg\_water * kg\_air\textasciicircum\{\}-1 & Water vapor contained in the outside air\\
\texttt{phi7} & 16.7 & m\textasciicircum\{\}3 * s\textasciicircum\{\}-1 & Capacity of air flow through the pad\\
\texttt{phi8} & 666.6 & m\textasciicircum\{\}3 * s\textasciicircum\{\}-1 & Air flow capacity of forced ventilation system\\
\texttt{phi9} & 0.916 & kg * s\textasciicircum\{\}-1 & Fog system capacity\\
\texttt{psi1} & 18 & kg * kmol\textasciicircum\{\}-1 & Molar mass of water\\
\texttt{psi2} & 13300 & mg * s\textasciicircum\{\}-1 & Capacity of the external CO2 source\\
\texttt{psi3} & 30.031 & g * mol\_CH2O\textasciicircum\{\}-1 & Molar mass of the CH2O\\
\texttt{rho1} & 0.07 &  & PAR reflection coefficient of the cannopy\\
\texttt{rho2} & 0.65 &  & Floor reflection coefficient PAR\\
\texttt{rho3} & 1.2 & kg * m\textasciicircum\{\}-3 & Air density\\
\texttt{rho4} & 1.2 & kg * m\textasciicircum\{\}-3 & Air density see level\\
\texttt{sigma} & 5.67e-08 & W * m\textasciicircum\{\}-2 * K\textasciicircum\{\}-4 & Stefan-Boltzmann constant\\
\texttt{tau1} & 0.77 &  & PAR transmission coefficient of the Cover\\
\texttt{tau2} & 1 &  & FIR transmission coefficient of the Cover\\
\texttt{tau3} & 0.11 &  & FIR transmission coefficient of the thermal screen\\
\bottomrule
\end{longtable}


\section{Photosynthesis constants}\label{app:photo}
Parameters of the FvCB leaf model and its temperature responses (Section~\ref{sec:photo}).
\begin{longtable}{@{}p{0.16\textwidth} p{0.16\textwidth} p{0.20\textwidth} p{0.42\textwidth}@{}}
\toprule
\textbf{Name} & \textbf{Value} & \textbf{Units} & \textbf{Description}\\
\midrule
\endfirsthead
\toprule
\textbf{Name} & \textbf{Value} & \textbf{Units} & \textbf{Description}\\
\midrule
\endhead
\texttt{ab} & 0.85 &  & Leaf absorbance\\
\texttt{alpha} & 0.35 &  & Product of the leaf absrovance, the fraction of light directed to PSII and the maximum quantum effiuciency of PSII; initial slope of an J/Iin curve\\
\texttt{C\_ev1} & 4.97685e-11 & mol\_phot * m\textasciicircum\{\}-2 * s\textasciicircum\{\}-1 & Constant in the formula of f\_R\\
\texttt{C\_ev2} & 6.25e-12 & mol\_phot * m\textasciicircum\{\}-2 * s\textasciicircum\{\}-1 & Constant in the formula of f\_R\\
\texttt{C\_ev3d} & 0.61 & mol\_air * mol\_CO2\textasciicircum\{\}-1 & Constant in the formula of f\_C\\
\texttt{C\_ev3n} & 1.1e-05 & mol\_air * mol\_CO2\textasciicircum\{\}-1 & Constant in the formula of f\_C\\
\texttt{C\_ev4d} & 4.3e-06 & Pa\textasciicircum\{\}-2 & Constant in the formula of f\_C\\
\texttt{C\_ev4n} & 5.2e-06 & Pa\textasciicircum\{\}-2 & Constant in the formula of f\_C\\
\texttt{E\_Soc} & -24460 & kg * m\textasciicircum\{\}2 * s\textasciicircum\{\}-2 * mol\textasciicircum\{\}-1 & Sco activation energy\\
\texttt{f} & 0.15 &  & Correction factor for the spectral quality of the light\\
\texttt{fc} & 2.95331e+10 & g * s * mol\_CO2\textasciicircum\{\}-1 & Factor to transform \(\mu\)mol\_CO2/sec to grams\_CH20/day\\
\texttt{J\_max} & 0.0002 & mol\_phot * m\textasciicircum\{\}-2 * s\textasciicircum\{\}-1 & Maximum electron transport rate\\
\texttt{k\_Ag} & 1e-06 & m\textasciicircum\{\}3 * g\textasciicircum\{\}-1 * s\textasciicircum\{\}-1 * mol\_CO2 * mol\_air\textasciicircum\{\}-1 & Constant for units transformation\\
\texttt{K\_C25} & 0.0003 & mol\_CO2 * mol\_air\textasciicircum\{\}-1 & Michaelis-Menten constant for CO2\\
\texttt{k\_d} & 86400 &  & Factor to transform s\textasciicircum\{\}-1 into d\textasciicircum\{\}-1\\
\texttt{k\_fc} & 1e-06 & mol\_CO2 * mol\_air\textasciicircum\{\}-1 & Constant for units completion\\
\texttt{k\_JV} & 1 &  & Auxiliary constant which transforms the units of the electron transport rate J to those of the maximum Rubisco rate V\_cmax\\
\texttt{K\_O25} & 0.3 & mol\_O2 * mol\_air\textasciicircum\{\}-1 & Michaelis-Menten constant for O2\\
\texttt{k\_T} & 1 & C & Auxiliary constant to add temperature units\\
\texttt{ks} & 0.5 &  & Stomatal ratio\\
\texttt{O\_a} & 0.21 & mol\_O2 * mol\_air\textasciicircum\{\}-1 & O2 concentration in the environment\\
\texttt{phi} & 2 & mol\_O2 * mol\_CO2\textasciicircum\{\}-1 & Ratio of oxygenation to carboxylation rates\\
\texttt{Q10\_KC} & 2.1 &  & Temperature response of Michaelis-Menten constant for CO2\\
\texttt{Q10\_KO} & 1.2 &  & Temperature response of Michaelis-Menten constant for O2\\
\texttt{Q10\_tau} & 2.1 &  & Temperature response of specificity factor\\
\texttt{Q10\_Vcmax} & 2.4 &  & V\_cmax value at 25C\\
\texttt{r\_m} & 100 & s * m\textasciicircum\{\}-1 & Minimal stomatal resistance\\
\texttt{Rb} & 711 & s * m\textasciicircum\{\}-1 & Stomatal resistance of the canopy boundary layer\\
\texttt{Rd\_day} & 1e-06 & mol * m\textasciicircum\{\}-2 * s\textasciicircum\{\}-1 & Rate of respiratory CO2 released from the mitocondria\\
\texttt{Rgas} & 8.31447 & m\textasciicircum\{\}2 * kg * s\textasciicircum\{\}-2 * K\textasciicircum\{\}-1 * mol\textasciicircum\{\}-1 & Universal Gas constant\\
\texttt{Rs} & 5e-06 & mol\_phot * m\textasciicircum\{\}-2 * s\textasciicircum\{\}-1 & Radiation setpoint to switch day and night\\
\texttt{S} & -86400 & m\textasciicircum\{\}2 * s * mol\_phot\textasciicircum\{\}-1 & Constant in the formula of Sr\\
\texttt{Sco25} & 2800 & Pa * Pa\textasciicircum\{\}-1 & Rubisco CO2/O2 specificity factor (Sc/o) at 25C\\
\texttt{tau\_25} & 2600 & mol * mol\textasciicircum\{\}-1 & Specificity factor\\
\texttt{theta} & 0.7 &  & Empirical curvature factor\\
\texttt{V\_cmax25} & 0.00015 & mol\_CO2 * m\textasciicircum\{\}-2 * s\textasciicircum\{\}-1 & Maximum Rubisco Rate, per unit area\\
\bottomrule
\end{longtable}


\section{Crop-growth constants}\label{app:growth}
Parameters of the cucumber source--sink growth model, including flowering, fruit and the leaf-cohort canopy (Sections~\ref{sec:growth}).
\begin{longtable}{@{}p{0.16\textwidth} p{0.16\textwidth} p{0.20\textwidth} p{0.42\textwidth}@{}}
\toprule
\textbf{Name} & \textbf{Value} & \textbf{Units} & \textbf{Description}\\
\midrule
\endfirsthead
\toprule
\textbf{Name} & \textbf{Value} & \textbf{Units} & \textbf{Description}\\
\midrule
\endhead

\bottomrule
\end{longtable}


\section{State variables, inputs and controls}\label{app:glossary}

\begin{longtable}{@{}p{0.12\textwidth} p{0.20\textwidth} p{0.62\textwidth}@{}}
\caption{Dynamic state variables (climate ODE).}\label{tab:state}\\
\toprule \textbf{Symbol} & \textbf{Units} & \textbf{Meaning}\\ \midrule \endfirsthead
\toprule \textbf{Symbol} & \textbf{Units} & \textbf{Meaning}\\ \midrule \endhead
$T_1$ & K & Canopy (leaf) temperature\\
$T_2$ & K & Greenhouse-air temperature (the controlled variable)\\
$C_1$ & mg m$^{-3}$ & Greenhouse-air \ce{CO2} concentration\\
$V_1$ & Pa & Greenhouse-air vapour pressure\\
$T_5$ & K & Floor / soil thermal-mass temperature\\
\bottomrule
\end{longtable}

\begin{longtable}{@{}p{0.10\textwidth} p{0.22\textwidth} p{0.62\textwidth}@{}}
\caption{Inputs to the climate right-hand side (weather drivers and controller-supplied quantities).}\label{tab:inputs}\\
\toprule \textbf{Symbol} & \textbf{Units} & \textbf{Meaning}\\ \midrule \endfirsthead
\toprule \textbf{Symbol} & \textbf{Units} & \textbf{Meaning}\\ \midrule \endhead
$I_1$ & m$^2$ m$^{-2}$ & Leaf area index (LAI), from the growth model\\
$I_2$ & W m$^{-2}$ & Outside global (short-wave) radiation\\
$I_3$ & K & Heating-pipe temperature (from the heating controller)\\
$I_4$ & K & Effective sky temperature (Berdahl--Martin)\\
$I_5$ & K & Outside air temperature\\
$I_6$ & K & Mechanical-cooling temperature\\
$I_7$ & K & Floor temperature seen by the FIR/conduction terms ($=T_5$)\\
$I_8$ & m s$^{-1}$ & Outside wind speed\\
$I_9$ & W m$^{-2}$ & PAR reaching the canopy $=(1-\eta_1)\tau_1\eta_2\,I_{docel}$\\
$I_{10}$ & mg m$^{-3}$ & Outside \ce{CO2} concentration\\
$I_{11}$ & Pa & Outside vapour pressure\\
$T_{deep}$ & K & Deep-soil (anchor) temperature, from weather\\
\bottomrule
\end{longtable}

\begin{longtable}{@{}p{0.10\textwidth} p{0.30\textwidth} p{0.54\textwidth}@{}}
\caption{Control actuators $U_1$--$U_{12}$ (each in $[0,1]$ unless noted).}\label{tab:controls}\\
\toprule \textbf{Symbol} & \textbf{Actuator} & \textbf{Role / entry point}\\ \midrule \endfirsthead
\toprule \textbf{Symbol} & \textbf{Actuator} & \textbf{Role / entry point}\\ \midrule \endhead
$U_1$ & Thermal / energy screen & FIR terms $r_{12}$, screen transmission $g_2,g_3$; deploy/crack\\
$U_2$ & Fan--pad (evaporative) inlet & $h_2,p_2,p_6,o_3$ (0 in reference config)\\
$U_3$ & Mechanical cooling & (0 in reference config)\\
$U_4$ & Direct air heater & $o_1,p_4$ (0 in reference config)\\
$U_5$ & Shade/black-out screen & modifies vent discharge $n_1,n_3$ (0 in ref.)\\
$U_6$ & Side (roll-up) ventilation & wind-driven exchange $f_5$\\
$U_7$ & Forced ventilation (fans) & $f_3$ (0 in reference config)\\
$U_8$ & Roof (ridge) ventilation & buoyancy+wind exchange $f_7$\\
$U_9$ & Fogging / misting & $p_3$ (latent) and $l_2$\\
$U_{10}$ & External \ce{CO2} dosing valve & source $o_2$\\
$U_{11}$ & (reserved / measurement) & --\\
$U_{12}$ & Supplemental lamps (on/off) & PAR to canopy; $r_5$ (NIR), $h_{12}$ (heat)\\
\bottomrule
\end{longtable}

\end{document}
